\chapter{The Optical Model}

\sourcepage{1}
\section{From a real well to absorption}
An explicit square-well model is
\[
  V(r)=
  \begin{cases}
    -V_0, & r\le R,\\
    0, & r>R,
  \end{cases}
\]
where $R$ is the nuclear radius.  The benefit is that it is very easy to use.  The issue is that a real potential cannot explain the low-energy capture observed by Szilard, Fermi, and others.

The next-best step is to enforce absorption of neutrons that are not reflected from the target, following Feshbach (1949).  In a one-dimensional planar analogue, take
\[
  \psi^{I}=\ee^{\ii kx}+A\ee^{-\ii kx},
  \qquad
  \psi^{II}=B\ee^{\ii\alpha x},
\]
with
\[
  k=\frac{\sqrt{2mE}}{\hbar},
  \qquad
  \alpha=\frac{\sqrt{2m(E+V_0)}}{\hbar}.
\]
Continuity of the wave function and derivative at the interface gives
\[
  1+A=B,
  \qquad
  \ii k-\ii kA=\ii\alpha B.
\]
Therefore
\[
  B=\frac{2k}{\alpha+k},
  \qquad
  A=\frac{k-\alpha}{k+\alpha}.
\]

\sourcepage{2}
The transmission coefficient is
\begin{align*}
  T
  &=\frac{|1|^2-|A|^2}{|1|^2}\\
  &=1-\left(\frac{\alpha-k}{\alpha+k}\right)^2\\
  &=\frac{4\alpha k}{(\alpha+k)^2}.
\end{align*}
A semiclassical capture cross section is then written
\[
  \sigma_c=\pi(R+\bar\lambda)^2T,
\]
which is the ``black nucleus'' model.  A typical well depth is $V_0\simeq40\unit{MeV}$, corresponding to $\alpha$ of order $10^{13}\unit{cm^{-1}}$.

The benefit is that this model captures the $1/v$ behavior.  The issue is that it produces a monotonic $\sigma(E)$ and therefore does not reproduce fine resonance structure.  The source includes the reminder: ``Add plot to demonstrate.''

\sourcepage{3}
The same idea is equivalent to an interior converging wave,
\[
  r\psi\sim\ee^{-\ii\alpha r},
  \qquad r<a,
\]
so that neutrons flow into the target without an outgoing interior component.  The corresponding logarithmic derivative is
\[
  D=a\frac{\dd(r\psi)/\dd r}{r\psi}=-\ii\alpha a.
\]
\editorial{A sticky note on the scan asks whether this should be $-\ii\alpha a$; the expression shown here follows the derivative of $\ee^{-\ii\alpha r}$.}

Use the same matching procedure as before.  For the exterior partial wave,
\[
  \chi_\ell^{\mathrm{ext}}(kr)
  =H_\ell^-(kr)-S_\ell H_\ell^+(kr)
  \sim \ee^{-\ii(kr-\ell\pi/2)}
  -S_\ell\ee^{\ii(kr-\ell\pi/2)}.
\]
Matching the logarithmic derivative at $r=a$ gives an $S$-matrix ratio of the form
\[
  S_\ell
  =\frac{H_\ell^{-\prime}-(D/a)H_\ell^-}
  {H_\ell^{+\prime}-(D/a)H_\ell^+}.
\]
For $\ell=0$, this reduces to the same reflection-amplitude behavior as the planar model,
\[
  S_0\sim\frac{\alpha-k}{\alpha+k},
\]
so $1-|S_0|^2$ gives the absorbed fraction used in the usual definitions of $\sigma_{r,\ell}$ and related cross sections.

\sourcepage{4}
\section{Feshbach's question and the strong-coupling model}
The main question is: does a particle necessarily form a compound nucleus after entering the target?  In other words, is the strong-coupling model valid?

With less interaction, assume that the target acts first as a potential well.  A particle has a chance to enter and, once inside, a chance to form a compound nucleus.  The notes organize the cross sections as
\begin{align*}
  \sigma_t&=\sigma_{\mathrm{el}}+\sigma_r,\\
  \sigma_{\mathrm{el}}&=\sigma_{\mathrm{pe}}+\sigma_{\mathrm{el},c},\\
  \sigma_c&=\sigma_{\mathrm{el},c}+\sigma_r,
\end{align*}
where ``pe'' denotes potential elastic scattering and the subscript $c$ denotes compound-nucleus contributions.

Model the target with a complex potential,
\[
  V(r)=
  \begin{cases}
    V_0+\ii V_1, & r<a,\\
    0, & r\ge a.
  \end{cases}
\]
Such a model should yield scattering that reproduces the potential-elastic component, while the absorptive part represents the compound-nucleus cross section averaged over many resonances.

\sourcepage{5}
The optical-model sketch separates three conceptual stages:
\begin{enumerate}
  \item A projectile enters the optical potential, where it may be elastically scattered or absorbed.
  \item Absorption creates a ``compound system,'' including interactions of the projectile with part of the target, such as the surface, or a collective excitation in which the nucleus is excited together.
  \item The compound system evolves and decays, possibly through a fully equilibrated compound nucleus.
\end{enumerate}
Thus the treatment moves from Bohr's two stages to Weisskopf's three stages.

The optical potential is non-Hermitian so that absorption is produced.  Its simplest form is
\[
  V(r)=
  \begin{cases}
    V_0+\ii V_1, & r<a,\\
    0, & r\ge a.
  \end{cases}
\]

\sourcepage{6}
\section{Energy averaging and the optical \texorpdfstring{$S$}{S} matrix}
Previously,
\begin{align*}
  \sigma_{\mathrm{el},\ell}
  &=\pi\bar\lambda^2(2\ell+1)|1-S_\ell|^2,\\
  \sigma_{r,\ell}
  &=\pi\bar\lambda^2(2\ell+1)(1-|S_\ell|^2),\\
  \sigma_{t,\ell}
  &=2\pi\bar\lambda^2(2\ell+1)(1-\operatorname{Re}S_\ell),
\end{align*}
with
\[
  S_\ell=\ee^{2\ii\delta_\ell}.
\]
Define an energy-averaged $S$ matrix by
\begin{equation}
  \brac{S_\ell(E)}
  \equiv\frac{1}{\Delta E}
  \int_{E-\Delta E/2}^{E+\Delta E/2}S_\ell(E')\dd E'.
  \label{eq:om-average-S}
\end{equation}
The averaged cross sections are
\begin{align*}
  \overline\sigma_{\mathrm{el},\ell}
  &=\pi\bar\lambda^2(2\ell+1)\brac{|1-S_\ell|^2},\\
  \overline\sigma_{r,\ell}
  &=\pi\bar\lambda^2(2\ell+1)\left[1-\brac{|S_\ell|^2}\right],\\
  \overline\sigma_{t,\ell}
  &=2\pi\bar\lambda^2(2\ell+1)
  \left[1-\operatorname{Re}\brac{S_\ell}\right].
\end{align*}
The averaging interval is assumed small enough that $\bar\lambda^2$ need not itself be averaged.

Now
\begin{align*}
  \brac{|1-S_\ell|^2}
  &=1-\brac{S_\ell}-\brac{S_\ell^*}+\brac{|S_\ell|^2}\\
  &=|1-\brac{S_\ell}|^2
  +\brac{|S_\ell|^2}-|\brac{S_\ell}|^2.
\end{align*}

\sourcepage{7}
This decomposition defines
\begin{equation}
  \overline\sigma_{\mathrm{el},p,\ell}
  =\pi\bar\lambda^2(2\ell+1)|1-\brac{S_\ell}|^2,
  \label{eq:om-potential-elastic}
\end{equation}
and
\begin{equation}
  \overline\sigma_{\mathrm{el},c,\ell}
  =\pi\bar\lambda^2(2\ell+1)
  \left(\brac{|S_\ell|^2}-|\brac{S_\ell}|^2\right).
  \label{eq:om-compound-elastic}
\end{equation}
The compound-formation cross section is then
\begin{align}
  \overline\sigma_{c,\ell}
  &=\overline\sigma_{r,\ell}
  +\overline\sigma_{\mathrm{el},c,\ell}\\
  &=\pi\bar\lambda^2(2\ell+1)
  \left(1-|\brac{S_\ell}|^2\right)\\
  &=\pi\bar\lambda^2(2\ell+1)T_\ell(E),
  \label{eq:om-transmission-coefficient}
\end{align}
where
\[
  T_\ell(E)=1-|\brac{S_\ell}|^2
\]
is the transmission coefficient.

Usually $\brac{S_\ell}$ yields the ``gross-structure'' problem.  The potential-elastic contribution is the pure one-particle scattering part.  At very large energy, $\brac{S_\ell}\to S_\ell$ and the compound-elastic fluctuation term vanishes.  In a compound-nucleus model, the probability for the system to decay back to the entrance channel is then small because many other exit channels are available.  The averaged elastic and compound-formation cross sections can therefore be described by a one-particle complex potential.

\sourcepage{8}
Both the black-nucleus and optical models yield
\begin{equation}
  \sigma_{c,\ell}
  =\pi\bar\lambda^2(2\ell+1)(1-|S_\ell|^2)
  =\pi\bar\lambda^2(2\ell+1)T_\ell(E),
\end{equation}
where $T_\ell(E)$ is the transmission coefficient.

The sketch of neutron cross section on $^{235}$U distinguishes low-energy behavior, isolated resonances with a background, and overlapping resonances beyond which the average optical description captures the physics.

A more complicated optical potential may be written schematically as
\[
  V(r)
  =(V_0+\ii V_1)f_{\mathrm{vol}}(r)
  +(V_2+\ii V_3)f_{\mathrm{surf}}(r)
  +\bm\ell\cdot\bm s\,(V_4+\ii V_5)f_{\mathrm{so}}(r),
\]
with volume, surface, and spin--orbit form factors.

\sourcepage{9}
\section{Hauser--Feshbach connection}
Bohr's hypothesis gives
\[
  \sigma_{\alpha\beta}
  =\sigma_c(\alpha)\frac{\Gamma_\beta}{\Gamma},
\]
where the branching ratio is the decay probability into channel $\beta$.

The optical model supplies the channel transmission coefficient,
\[
  \sigma_{c,\ell}^{(\alpha)}
  =\pi\bar\lambda_\alpha^2(2\ell+1)T_{\alpha\ell},
  \qquad
  T_{\alpha\ell}=1-|S_{\alpha\ell}|^2.
\]
Reciprocity or detailed balance gives
\[
  k_\alpha^2\sigma_{\alpha\beta}
  =k_\beta^2\sigma_{\beta\alpha}.
\]
Neglecting spins, the basic Hauser--Feshbach form is
\begin{equation}
  \sigma_{\alpha\beta}
  =\pi\bar\lambda_\alpha^2(2\ell+1)
  \frac{T_\alpha T_\beta}{\sum_\gamma T_\gamma}.
  \label{eq:om-hauser-feshbach}
\end{equation}
The expression becomes much more elaborate when spins and angular dependence are included, but the basic structure is unchanged.

In the Weisskopf--Ewing approximation, the channel transmissions are replaced or supplemented by a continuum level density for the residual nucleus.  The source points to the Wigner distribution as the next example.
